% MathSearch — guide and test sheet.
%
% The document the viewer opens by default (lib/pdfjs/web/mathsearch-guide.pdf,
% built by `npm run guide`). It explains the extension and doubles as a test
% sheet: every formula and statement in it is made up for this purpose, so
% it can be shipped with the extension.
%
% It is laid out to exercise the features:
%   - roman-numbered front matter, so printed page 1 is PDF page 3
%     ("Look up in index" must map printed -> PDF pages);
%   - hyperref bookmarks and page labels;
%   - an index of notation at the end whose page numbers come from \pageref,
%     so they stay correct when the text changes;
%   - lookalike symbols (\mathfrak{F}, \mathscr{F}, \mathcal{F}, F) and
%     symbols with and without sub/superscripts.
\documentclass[11pt]{article}
\usepackage[margin=1in]{geometry}
\usepackage{amsmath,amssymb,mathrsfs}
\usepackage[hidelinks]{hyperref}

\newcommand{\key}[1]{\fbox{\small\ttfamily #1}}
\newcommand{\cmd}[1]{\texttt{\textbackslash #1}}

\begin{document}

% ── Front matter (roman page numbers) ───────────────────────────────────────
\pagenumbering{roman}
\begin{center}
  \vspace*{2cm}
  {\LARGE\bfseries MathSearch}\\[4pt]
  {\large Search a PDF for mathematical symbols by how they look}\\[18pt]
  {\large Guide and test sheet}\\[6pt]
  Version 1.0
\end{center}
\vspace{1.5cm}
\noindent This document opens when the MathSearch viewer starts. It explains
how to use the extension and is also a practice sheet: the formulas in it
exist only so that there is something to search for. They are made up and
state nothing of mathematical interest.

\medskip\noindent To search your own PDF, open it with the viewer's
\textbf{Open} button (in the \textbf{Tools} menu at the top right) or press
\key{Ctrl+O}, or drag the file onto the viewer. To read several PDFs at once, click the MathSearch toolbar
icon again: each click opens a new viewer tab with its own search.
\newpage
\tableofcontents
\newpage

% ── Main text (arabic page numbers) ─────────────────────────────────────────
\pagenumbering{arabic}

\section{Getting started}

The MathSearch panel floats over the page, at the top right when the viewer
opens. Drag it by its header to move it out of the way, and double-click the
header to put it back. The \textbf{$\blacktriangledown$} button collapses it
and \textbf{$\times$} hides it. The \textbf{$\Sigma$} button in the viewer's
toolbar, or \key{Alt+M}, brings it back.

\begin{enumerate}
  \item Type the symbol as LaTeX in the box, for example
        \verb|\mathfrak{F}|. The preview underneath shows it as it will be
        searched for: what you see is the search template.
  \item Press \key{Enter} or click \textbf{Search}. Every place in the
        document where the symbol \emph{looks} the same is highlighted in
        yellow, and the current match in orange, each with its score.
  \item Press \key{Enter} again to go to the next match and
        \key{Shift+Enter} to go back, or click an entry in the list.
\end{enumerate}

\subsection*{Searching by snip}
If you don't know a symbol's LaTeX, click the scissors button next to the
box (or press \key{Alt+S}) and drag a rectangle around the symbol on the
page. The snip replaces the preview and is searched for straight away. It
stays the query until you type in the box again. Snip only the symbol:
anything else inside the rectangle, such as a subscript, a fraction bar or
the edge of a neighbouring letter, is searched for too and lowers the scores
of clean matches. \key{Esc} cancels a snip.

\subsection*{Keyboard shortcuts}
\begin{center}
\begin{tabular}{ll}
  \key{Alt+M}        & show the panel and put the cursor in the search box \\
  \key{Enter}        & search; once there are results, next match \\
  \key{Shift+Enter}  & previous match \\
  \key{Ctrl+Enter}   & look the symbol up in the book's index first \\
  \key{Alt+S}        & snip a symbol from the page and search for it \\
  \key{Esc}          & leave the search box \\
  double-click header & move the panel back to the top right \\
\end{tabular}
\end{center}

\subsection*{Settings in the panel}
\textbf{Scope} searches the whole document or only the current page.
\textbf{Min score} (0.90 by default) hides weaker matches. Lowering it after
a search shows more candidates at once, without searching again. This helps
with small occurrences, such as a symbol inside a subscript.
\textbf{Penalise scripts and accents next to the match} (on by default)
makes a symbol that carries an extra sub- or superscript, accent or dot rank
below the bare symbol (see Section~\ref{sec:search}). Untick it to find the
symbol with or without them; the search runs again, because the scores
change. The setting lasts until you close the viewer.

\subsection*{Memory settings}
Open \textbf{Settings} at the bottom of the panel. These two are kept in your
browser between sessions; nothing else is stored.
\textbf{Memory for the page index} (250\,MB by default, at least 50) limits
how much memory the extension keeps for pages it has already prepared for
searching. A page of text takes about 0.16\,MB, so most books fit. In a
longer book, the extension keeps the book's index and list of notation and
the pages around where your last search started, half before and half
after; other pages are prepared again when a search reaches them, which is
slower but finds the same matches. The line under the setting shows how much
is in use.
\textbf{Release it when the tab is hidden for} (30\,minutes by default;
also 10\,minutes, 1 or 2 hours, or never) gives the memory back while you
work in another tab. It is prepared again when you come back.

\section{How the search works}\label{sec:search}

MathSearch does not read the text of the PDF: it compares shapes. Each page
is rendered as an image, and the symbol from the preview is looked for in
it, at any size. This finds symbols that have no text form in the PDF, and
it tells apart letters that text search treats as the same, such as
$\mathfrak{F}$, $\mathscr{F}$, $\mathcal{F}$ and $F$.

\medskip\noindent\textbf{Search order.} The search starts at the page you are
reading and goes backward to the first page, then forward from where you
started to the end. A symbol is usually defined before it is used, so the
nearest earlier occurrence, often its definition, is found first. The viewer
jumps to the first match as soon as it is found, and the rest appear while
the search continues. \textbf{Cancel} stops it and keeps what was found.

\medskip\noindent\textbf{How long it takes.} After a book is opened, its pages
are prepared in the background (``Indexing pages 3/24 \dots''); a search
does not wait for that, but goes faster once it is done. A search through a
whole prepared book takes a few hundredths of a second per page, longer for
larger or compound symbols: in a 1100-page textbook on a laptop, about 40\,s
for $\mathscr{X}$ and about 90\,s for a snip of $L(H)$. Pages beyond the
memory limit (see \emph{Memory settings}) are prepared again during the
search, which adds to that: the same $\mathscr{X}$ search took about 70\,s
with the limit set to 50\,MB. Because the search starts at your page, the
nearest occurrence usually shows up within seconds, long before the search
ends.

\medskip\noindent\textbf{Sub- and superscripts count.} Searching for
$\mathbb{R}$ ranks $\mathbb{R}^d$ and $\hat{R}$ below a bare $\mathbb{R}$,
because the extra marks on the page are not part of what you asked for.
To find $\mathbb{R}^d$ itself, type \verb|\mathbb{R}^d|; to find
$\mathbb{R}$ with or without marks, untick \textbf{Penalise scripts and
accents}. Punctuation is not a script: in ``take $x$, then'' the comma sits
where a subscript would, but it is recognised by its shape and $x$ scores as
if it stood alone.

\section{Looking a symbol up in the index}\label{sec:lookup}

Many books have an index of notation that says where each symbol is
introduced. \textbf{Look up in index} (or \key{Ctrl+Enter}) uses it:
\begin{enumerate}
  \item It finds the index through the PDF's bookmarks, or by its heading:
        a list of notation or symbols near the start of the book (for
        example right after the contents), an index on one of the last
        pages, or both.
  \item It searches the index for the symbol and reads the page number
        printed next to it.
  \item It jumps to that page and searches from there.
\end{enumerate}
The printed page number is usually not the PDF page number, because front
matter comes first. In this document, printed page~1 is PDF page~3. The
panel shows where the index pointed, for example
``Index $\to$ p.\,3''. An ``$\approx$'' after a page means the PDF page
could not be confirmed and is an estimate.

This document has an index of notation on page~\pageref{sec:notation}. Try
it: type \verb|\mathfrak{F}| and press \key{Ctrl+Enter}.

\section{Practice text}\label{sec:practice}

The following text is made up. It uses each symbol in the index of notation
several times, in the ways they appear in real papers.

\subsection*{A toy filtration}
\textbf{Definition 1.}\label{def:frak} A \emph{toy filtration} is a family
$\mathfrak{F} = (\mathfrak{F}_n)_{n \ge 0}$ of finite sets with
$\mathfrak{F}_n \subseteq \mathfrak{F}_{n+1}$ for every $n$. We write
$\mathscr{F}$ for the union of all $\mathfrak{F}_n$.\label{def:scr}

\medskip\noindent\textbf{Definition 2.}\label{def:cal} For a toy filtration
$\mathfrak{F}$, let $\mathcal{F}$ be the collection of subsets of
$\mathscr{F}$ that meet only finitely many of the sets
$\mathfrak{F}_{n+1} \setminus \mathfrak{F}_n$.

\medskip\noindent\textbf{Definition 3.}\label{def:tau} A \emph{stopping
index} is a sequence $\tau_1 \le \tau_2 \le \cdots$ of natural numbers. We
write $\tau_j$ for its $j$-th term and $\mathscr{F}_{\tau_{j+1}}$ for the
set $\mathfrak{F}_{\tau_{j+1}}$ seen as a subset of $\mathscr{F}$.

\subsection*{Lookalikes}
The four letters $\mathfrak{F}$, $\mathscr{F}$, $\mathcal{F}$ and $F$ are
different symbols, and a search for one should not report the others.
The same goes for $\mathbb{R}$ and $R$.

\newpage
\subsection*{Spaces and classes}
\textbf{Definition 4.}\label{def:R} We write $\mathbb{R}$ for the real line
and $\mathbb{R}^d$ for $d$-dimensional space. A map
$f \colon \mathbb{R}^d \to \mathbb{R}$ is \emph{toy-smooth} if it is
constant on each set of $\mathcal{F}$.

\medskip\noindent\textbf{Definition 5.}\label{def:ID} Let
$\mathrm{ID}(\mathbb{R}^d)$ be the class of toy-smooth maps, and
$\mathrm{ID}_2(\mathbb{R}^d)$ those that also vanish outside a bounded set.

\medskip\noindent\textbf{Remark.} Small occurrences are harder to find:
in $\mathscr{F}_{\tau_j}$ the $\tau$ is in script size, and in
$e^{\mathfrak{F}}$ the $\mathfrak{F}$ is. Lower \textbf{Min score} to about
0.85 to see them.

\subsection*{Exercises}
Search for each of the following. The first page each one appears on is
given, so you can check that the search order and the index lookup take you
there.
\begin{center}
\begin{tabular}{lll}
  type & first defined & notes \\ \hline
  \verb|\mathfrak{F}| & p.\,\pageref{def:frak} & also in Section~\ref{sec:search} \\
  \verb|\mathscr{F}|  & p.\,\pageref{def:scr}  & should not match $\mathfrak{F}$ or $\mathcal{F}$ \\
  \verb|\mathcal{F}|  & p.\,\pageref{def:cal}  & \\
  \verb|\mathbb{R}^d| & p.\,\pageref{def:R}    & a bare $\mathbb{R}$ ranks lower \\
  \verb|\mathrm{ID}|  & p.\,\pageref{def:ID}   & also matches ``ID'' in IDEA, see Limitations \\
  \verb|\tau_j|       & p.\,\pageref{def:tau}  & also in script size \\
\end{tabular}
\end{center}

\section{Limitations}
\begin{itemize}
  \item Scanned books work for the visual search. Their page numbers
        cannot be read for \textbf{Look up in index}, because a scan has no
        text layer.
  \item Content printed sideways (a table turned on its side) is found
        when the PDF has a text layer: the search then also tries the symbol
        turned the same way. In a scanned book it is not found. Rotating the
        view in the viewer is fine.
  \item Very small symbols (below about 6 pixels high at the search
        resolution) are skipped: at that size all shapes look alike.
  \item Everything happens in your browser: no page or query leaves your
        computer.
  \item A symbol that looks the same as ordinary letters is found in running 
      text too: $\mathrm{ID}$ also matches the ``ID'' in IDEA, because the
      search compares shapes only.
  \item Commas and semicolons are told apart from subscripts by their shape,
      tuned on the font of this guide. In another font, a comma right after
      a symbol, as in ``$x$, $y$'', may still be taken for a subscript and
      push the symbol below \textbf{Min score}. If a symbol you can see is
      not found, untick \textbf{Penalise scripts and accents}.
\end{itemize}

% ── Index of notation ───────────────────────────────────────────────────────
\newpage
\phantomsection\label{sec:notation}
\addcontentsline{toc}{section}{Index of Notation}
\section*{Index of Notation}
\noindent
\begin{tabular}{@{}p{0.44\textwidth}@{\hspace{0.12\textwidth}}p{0.44\textwidth}@{}}
$\mathfrak{F}$ \dotfill toy filtration, \pageref{def:frak} &
$\mathbb{R}$ \dotfill real line, \pageref{def:R} \\
$\mathscr{F}$ \dotfill union of a filtration, \pageref{def:scr} &
$\mathbb{R}^d$ \dotfill $d$-dimensional space, \pageref{def:R} \\
$\mathcal{F}$ \dotfill finite-meeting sets, \pageref{def:cal} &
$\mathrm{ID}$ \dotfill toy-smooth maps, \pageref{def:ID} \\
$\tau_j$ \dotfill stopping index, \pageref{def:tau} &
$\mathrm{ID}_2$ \dotfill bounded toy-smooth maps, \pageref{def:ID} \\
\end{tabular}

\end{document}
